\documentclass[11pt]{article}
\usepackage{amsmath}
\title{Progress: HardcoreBosonChainLoss.jl}
\date{1 October 2026}
\begin{document}\maketitle
\section*{Scope}
The model is a one-dimensional interacting
hard-core boson chain with local loss, solved by MPS quantum-jump
trajectories. This package has Julia dependency metadata, a dense
small-chain Lindblad reference, tests, example, benchmark, README,
executed notebook, and model note. No GitHub package was reused because
the local ITensorMPS implementation and independent dense reference were
sufficient.
\section*{Commands and results}
With Julia 1.13.1, a writable temporary depot followed by the local
Julia depot, offline package mode, and single-threaded BLAS, the command
\begin{verbatim}
JULIA_DEPOT_PATH=/private/tmp/zeno-julia-depot:/Users/alexioschristopoulos/.julia JULIA_PKG_OFFLINE=true OPENBLAS_NUM_THREADS=1 /Users/alexioschristopoulos/.juliaup/bin/julia --project=. test/runtests.jl
\end{verbatim}
passed 19 of 19 checks in two test sets. The one-site dense occupation at
$\gamma=0.8$, $t=0.7$ was $0.571209064$, equal to $e^{-0.56}$ to nine
printed decimals. For a closed three-site interacting chain, the MPS
occupation error against dense unitary evolution was
$4.024\times10^{-6}$ at step $0.01$; the sum of occupations was
$2.000000000$. Reducing the step from $0.05$ to $0.01$ reduced the
error. A lossy three-site 400-record ensemble agreed with dense
Lindblad site occupations within its four-standard-error plus numerical
tolerance. The dense evolved density matrix had trace one and
nonnegative eigenvalues to roundoff.

Both examples ran. At four sites, $J=1$, $\gamma_2=0.5$, $t=0.6$,
and 400 records, the maximum site-occupation differences from dense
evolution were $0.0031$ for $V=0$ and $0.0029$ for $V=1$.
The benchmark reported $5.080$ s for 100 six-site trajectories and
20 steps after warm-up, with final mean total occupation $2.76000$.
This local timing is illustrative.

Four tutorial Julia cells executed sequentially in one process, saved
their outputs, and passed notebook-format validation. The notebook's
400-record lossy estimate at site 2 was $0.817288\pm0.013651$ versus
the dense value $0.813778$; its density-matrix trace printed as
$1.000000000$. Both LaTeX notes compiled with the desktop compiler.
No commit, push, publication, or release was performed.
\section*{Limitations}
Finite-step jump placement, MPS truncation, and Monte Carlo uncertainty
must each be controlled. Dense reference is capped at four sites.
Large entanglement can make the MPS trajectory representation costly.
The rates are prescribed Markovian losses rather than a microscopic bath.
\section*{Public-repository preparation, 2 October 2026 at 12:15 UTC}
Reviewed the owner's existing GitHub repositories, including
\texttt{Open-Systems-and-Current-Fluctuations} and
\texttt{Local\_Temperature}. Their vectorized-density-operator routines
do not simplify this pure-state MPS quantum-jump solver, so no code was
copied. Added comments and a README table connecting numerical quantities
to occupation, hopping, interaction, loss events, and Lindblad averages;
expanded the executed notebook's physical explanations. Added a Julia
\texttt{.gitignore} for generated manifests and notebook checkpoints.
The public GitHub repository was created by the owner.

With Julia 1.13.1, cached ITensors 0.9.14 and ITensorMPS 0.3.23, and one
BLAS thread, offline dependency instantiation succeeded. The documented
\texttt{test/runtests.jl} command passed 19 of 19 assertions. The example's
four-site maximum occupation differences from dense Lindblad evolution
were 0.0031 for $V=0$ and 0.0029 for $V=1$. The warmed six-site benchmark
took 4.111 s for 100 trajectories and 20 steps, with final mean total
occupation 2.76000. These local timings do not establish scaling.
\section*{Publication, 2 October 2026}
The owner-created public repository at
\texttt{github.com/AlexiosChristopoulos/Hardcore\_boson\_chain\_loss}
contains the ten package files, including the Julia \texttt{.gitignore},
source, README, example, benchmark, tests, LaTeX notes, and executed
notebook. The published \texttt{main} tree and README were fetched from
GitHub to confirm the upload. Generated \texttt{Manifest.toml} was
excluded. No code from the owner's other repositories was copied.
\end{document}
